\documentclass[12pt]{article}
\usepackage[utf8]{inputenc}
\usepackage[T1]{fontenc}
\usepackage{hyperref}
\title{Sandawe Language Resources for K-12 Curriculum Development}
\author{Compiled by Claude Sonnet 5.5 (claude-sonnet-5-5) for Andr\'as Kornai}
\date{}
\begin{document}
\maketitle

\section{Introduction}
Sandawe (ISO 639-3 sad) is a language isolate spoken in the Dodoma Region of Tanzania. The search found scholarly grammars and dictionaries but no school materials.

\section{Primary Educational Resources}
No primers, textbooks or school programmes for Sandawe turned up in the search.

\section{Linguistic and Reference Materials}
\begin{itemize}
\item Sander Steeman, \emph{A Grammar of Sandawe: A Khoisan Language of Tanzania} (LOT, 2012). Catalogue record: \url{https://ibali.uct.ac.za/s/ccse/item/17594}
\item Eric ten Raa (ed. C.\ and P.\ Ehret), \emph{A Dictionary of Sandawe} (R\"udiger K\"oppe Verlag, 2012): \url{https://www.koeppe.de/titel_a-dictionary-of-sandawe}
\item Ryohei Kagaya, \emph{A Classified Vocabulary of the Sandawe Language} (ILCAA, Tokyo, 1993). I found it only in a search summary and have no URL for it.
\item Starling database Sandawe PDF (comparative word list; content not checked): \url{https://starlingdb.org/new100/snd.pdf}
\end{itemize}

\section{Religious and Community Resources}
None found in the search.

\section{Digital and Audio Resources}
None found in the search.

\section{Curriculum Development Context}
Sandawe is documented well enough to support materials development, since a modern grammar and two dictionaries exist. No evidence of school use or a standard orthography appeared. Swahili is the language of schooling in Tanzania.
\end{document}